%%%%%%%%%%%%%%%%%%%%%%%%%%%%%%%%%%%%%%%%%%%%%%%%%%%%%%%%%%%%%%%%%%%
% Physik IV - Aufgabenblatt 4
% Standalone homework / learning / blackboard-talk version
%%%%%%%%%%%%%%%%%%%%%%%%%%%%%%%%%%%%%%%%%%%%%%%%%%%%%%%%%%%%%%%%%%%

\documentclass[11pt,a4paper]{scrartcl}

%%%%%%%%%%%%%%%%%%%%%%%%%%%%%%%%%%%%%%%%%%%%%%%%%%%%%%%%%%%%%%%%%%%
% Embedded file: includes/university_metadata.tex
%%%%%%%%%%%%%%%%%%%%%%%%%%%%%%%%%%%%%%%%%%%%%%%%%%%%%%%%%%%%%%%%%%%
\newcommand{\CourseTitle}{Physik IV: Kerne und Teilchen}
\newcommand{\CourseShortTitle}{Physik IV}
\newcommand{\DocumentKind}{Hausaufgabe, Tafelvortrag und Lernfassung}
\newcommand{\DocumentTitle}{Aufgabenblatt 4}
\newcommand{\DocumentSubtitle}{Bethe-Weizsaecker-Formel, Coulombterm, Volumenterm und Isobarenparabeln}
\newcommand{\AuthorName}{Tobias Laser}
\newcommand{\InstitutionName}{Universitaet Innsbruck}
\newcommand{\SubmissionDate}{\today}

\newcommand{\makeuniversitytitle}{%
    \begin{titlepage}
        \centering
        \vspace*{2cm}
        {\Large \CourseTitle\par}
        \vspace{1.2cm}
        {\huge\bfseries \DocumentTitle\par}
        \ifx\DocumentSubtitle\empty\else
            \vspace{0.4cm}
            {\Large \DocumentSubtitle\par}
        \fi
        \vspace{1cm}
        {\Large \DocumentKind\par}
        \vfill
        {\large \AuthorName\par}
        \vspace{0.2cm}
        {\large \InstitutionName\par}
        \vspace{0.8cm}
        {\large \SubmissionDate\par}
    \end{titlepage}
    \pagenumbering{arabic}
}

\newcommand{\makechaptertitle}{%
    \begin{center}
        {\Large \CourseTitle\par}
        \vspace{0.5cm}
        {\huge\bfseries \DocumentTitle\par}
        \ifx\DocumentSubtitle\empty\else
            \vspace{0.25cm}
            {\Large \DocumentSubtitle\par}
        \fi
        \vspace{0.5cm}
        {\large \AuthorName \hfill \SubmissionDate\par}
    \end{center}
    \vspace{0.5cm}
}
\newcommand{\makephysiktitle}{\makeuniversitytitle}

%%%%%%%%%%%%%%%%%%%%%%%%%%%%%%%%%%%%%%%%%%%%%%%%%%%%%%%%%%%%%%%%%%%
% Embedded file: includes/university_packages.tex
%%%%%%%%%%%%%%%%%%%%%%%%%%%%%%%%%%%%%%%%%%%%%%%%%%%%%%%%%%%%%%%%%%%
\usepackage[margin=2.5cm]{geometry}
\usepackage[onehalfspacing]{setspace}
\usepackage{scrlayer-scrpage}

\usepackage[T1]{fontenc}
\usepackage[utf8]{inputenc}
\usepackage[ngerman]{babel}
\usepackage{lmodern}
\usepackage{microtype}
\usepackage{csquotes}

\usepackage{amsmath,amsthm,amssymb,mathtools}
\usepackage{bm}
\usepackage{icomma}
\usepackage{cancel}

\usepackage{graphicx}
\usepackage{tikz}
\usetikzlibrary{arrows.meta,calc,decorations.pathmorphing,positioning,patterns,angles,quotes}
\usepackage{pgfplots}
\pgfplotsset{compat=1.18}

\usepackage[locale=DE,space-before-unit=true,per-mode=symbol,separate-uncertainty=true]{siunitx}
\usepackage{booktabs,multirow,array,tabularx,longtable}
\usepackage{caption,subcaption}
\usepackage{placeins}
\usepackage{enumitem}
\usepackage{xparse}

\usepackage[most]{tcolorbox}
\tcbuselibrary{breakable,skins,theorems}

\usepackage[breaklinks=true,colorlinks=true,linkcolor=blue,urlcolor=blue,citecolor=blue]{hyperref}
\usepackage[nameinlink,noabbrev]{cleveref}

\clearpairofpagestyles
\ihead{\CourseShortTitle}
\ohead{\DocumentKind}
\cfoot{\pagemark}
\pagestyle{scrheadings}

\setlist[enumerate]{itemsep=0.4em,topsep=0.4em}
\setlist[itemize]{itemsep=0.25em,topsep=0.35em}
\captionsetup{font=small,labelfont=bf}

%%%%%%%%%%%%%%%%%%%%%%%%%%%%%%%%%%%%%%%%%%%%%%%%%%%%%%%%%%%%%%%%%%%
% Embedded file: includes/university_macros.tex
%%%%%%%%%%%%%%%%%%%%%%%%%%%%%%%%%%%%%%%%%%%%%%%%%%%%%%%%%%%%%%%%%%%
\newcommand{\R}{\mathbb{R}}
\newcommand{\N}{\mathbb{N}}
\newcommand{\Z}{\mathbb{Z}}
\newcommand{\Q}{\mathbb{Q}}
\newcommand{\C}{\mathbb{C}}

\newcommand{\set}[1]{\left\{ #1 \right\}}
\newcommand{\abs}[1]{\left| #1 \right|}
\newcommand{\norm}[1]{\left\lVert #1 \right\rVert}
\newcommand{\avg}[1]{\left\langle #1 \right\rangle}
\newcommand{\paren}[1]{\left( #1 \right)}
\newcommand{\bracket}[1]{\left[ #1 \right]}
\newcommand{\ol}[1]{\overline{#1}}

\newcommand{\dd}{\,\mathrm{d}}
\newcommand{\dv}[2]{\frac{\dd #1}{\dd #2}}
\newcommand{\pdv}[2]{\frac{\partial #1}{\partial #2}}
\newcommand{\pddv}[2]{\frac{\partial^2 #1}{\partial #2^2}}
\newcommand{\evalat}[2]{\left. #1 \right|_{#2}}

\newcommand{\vect}[1]{\bm{#1}}
\newcommand{\unitvect}[1]{\hat{\bm{#1}}}
\newcommand{\grad}{\vect{\nabla}}
\newcommand{\divergence}{\grad\!\cdot}
\newcommand{\curl}{\grad\!\times}

\newcommand{\half}{\frac{1}{2}}
\newcommand{\third}{\frac{1}{3}}
\newcommand{\twothirds}{\frac{2}{3}}
\newcommand{\hbarc}{\hbar c}
\newcommand{\alphaf}{\alpha}
\newcommand{\eV}{\si{eV}}
\newcommand{\MeV}{\si{MeV}}
\newcommand{\GeV}{\si{GeV}}
\newcommand{\TeV}{\si{TeV}}

\newcommand{\nuclide}[3]{^{#1}_{#2}\mathrm{#3}}
\newcommand{\isotope}[2]{^{#1}\mathrm{#2}}
\newcommand{\anti}[1]{\overline{#1}}
\newcommand{\pbar}{\anti{p}}
\newcommand{\nbar}{\anti{n}}
\newcommand{\ee}{e^-}
\newcommand{\ep}{e^+}
\newcommand{\mup}{\mu^+}
\newcommand{\mum}{\mu^-}
\newcommand{\nue}{\nu_e}
\newcommand{\numu}{\nu_\mu}
\newcommand{\nutau}{\nu_\tau}

\newcommand{\diffxs}{\frac{\dd\sigma}{\dd\Omega}}
\newcommand{\totalxs}{\sigma_{\mathrm{total}}}
\newcommand{\lum}{\mathcal{L}}
\newcommand{\smand}{s}
\newcommand{\tmand}{t}
\newcommand{\umand}{u}
\newcommand{\Ecm}{E_{\mathrm{CM}}}

\newcommand{\ie}{d.\,h.}
\newcommand{\eg}{z.\,B.}
\newcommand{\wrt}{bezüglich}

\DeclareSIUnit{\barn}{b}
\DeclareSIUnit{\millibarn}{mb}
\DeclareSIUnit{\microbarn}{\micro b}
\DeclareSIUnit{\nanobarn}{nb}
\DeclareSIUnit{\fermi}{fm}
\DeclareSIUnit{\year}{a}

%%%%%%%%%%%%%%%%%%%%%%%%%%%%%%%%%%%%%%%%%%%%%%%%%%%%%%%%%%%%%%%%%%%
% Embedded file: includes/university_boxes.tex
%%%%%%%%%%%%%%%%%%%%%%%%%%%%%%%%%%%%%%%%%%%%%%%%%%%%%%%%%%%%%%%%%%%
\tcbset{
    unibox/.style={
        enhanced,
        breakable,
        sharp corners,
        boxrule=0.6pt,
        left=1.2mm,
        right=1.2mm,
        top=1mm,
        bottom=1mm,
        before skip=0.8em,
        after skip=0.8em,
        fonttitle=\bfseries
    }
}

\newtcolorbox{calculationbox}[1][]{unibox,title=Rechnung,colback=black!2,colframe=black!55,#1}
\newtcolorbox{sidecalcbox}[1][]{unibox,title=Nebenrechnung,colback=black!1,colframe=black!40,#1}
\newtcolorbox{solutionbox}[1][]{unibox,title=Loesung,colback=blue!2,colframe=blue!45!black,#1}
\newtcolorbox{answerbox}[1][]{unibox,title=Antwort,colback=green!3,colframe=green!45!black,#1}
\newtcolorbox{notebox}[1][]{unibox,title=Notiz,colback=yellow!6,colframe=yellow!45!black,#1}
\newtcolorbox{definitionbox}[1][]{unibox,title=Definition,colback=cyan!3,colframe=cyan!45!black,#1}
\newtcolorbox{warningbox}[1][]{unibox,title=Achtung,colback=red!3,colframe=red!55!black,#1}
\newtcolorbox{summarybox}[1][]{unibox,title=Zusammenfassung,colback=violet!3,colframe=violet!50!black,#1}
\newtcolorbox{formulabox}[1][]{unibox,title=Formel,colback=orange!4,colframe=orange!65!black,#1}
\newtcolorbox{taskbox}[1][]{unibox,title=Aufgabe,colback=white,colframe=black!75,#1}

%%%%%%%%%%%%%%%%%%%%%%%%%%%%%%%%%%%%%%%%%%%%%%%%%%%%%%%%%%%%%%%%%%%
% Embedded file: includes/university_theorems.tex
%%%%%%%%%%%%%%%%%%%%%%%%%%%%%%%%%%%%%%%%%%%%%%%%%%%%%%%%%%%%%%%%%%%
\theoremstyle{definition}
\newtheorem{definition}{Definition}[section]
\newtheorem{example}{Beispiel}[section]
\theoremstyle{plain}
\newtheorem{satz}{Satz}[section]
\newtheorem{lemma}{Lemma}[section]
\theoremstyle{remark}
\newtheorem*{remark}{Bemerkung}

%%%%%%%%%%%%%%%%%%%%%%%%%%%%%%%%%%%%%%%%%%%%%%%%%%%%%%%%%%%%%%%%%%%
% Local macros for this sheet
%%%%%%%%%%%%%%%%%%%%%%%%%%%%%%%%%%%%%%%%%%%%%%%%%%%%%%%%%%%%%%%%%%%
\newcommand{\Eab}{E_{\mathrm{ab}}}
\newcommand{\EWW}{E_{\mathrm{WW}}}
\newcommand{\aV}{a_{\mathrm V}}
\newcommand{\aS}{a_{\mathrm S}}
\newcommand{\aC}{a_{\mathrm C}}
\newcommand{\aA}{a_{\mathrm A}}
\newcommand{\aP}{a_{\mathrm P}}
\newcommand{\mpmass}{m_{\mathrm p}}
\newcommand{\mnmass}{m_{\mathrm n}}
\newcommand{\memass}{m_{\mathrm e}}

\begin{document}

\makeuniversitytitle
\tableofcontents
\newpage

%%%%%%%%%%%%%%%%%%%%%%%%%%%%%%%%%%%%%%%%%%%%%%%%%%%%%%%%%%%%%%%%%%%
\section*{Kurzueberblick fuer das Lernen}
\addcontentsline{toc}{section}{Kurzueberblick fuer das Lernen}

\begin{summarybox}[title=Was dieses Blatt prueft]
Dieses Blatt verbindet die Bethe-Weizsaecker-Massenformel mit einfachen physikalischen Modellen:
\begin{itemize}
    \item Der \textbf{Coulombterm} entsteht aus der elektrostatischen Selbstenergie einer homogen geladenen Kugel.
    \item Der \textbf{Volumenterm} entsteht modellhaft aus der kurzreichweitigen Bindung zwischen naechsten Nachbarn.
    \item Die \textbf{Isobarenparabel} entsteht, weil die Masse bei festem $A$ als Funktion von $Z$ naeherungsweise quadratisch ist.
\end{itemize}
Der wichtigste Lernpunkt ist nicht nur das Einsetzen von Zahlen, sondern die Interpretation: einfache Kernmodelle liefern oft die richtige Groessenordnung, aber nicht alle mikroskopischen Details.
\end{summarybox}

\begin{formulabox}[title=Bethe-Weizsaecker-Massenformel in der verwendeten Form]
\[
M(A,Z)=(A-Z)m_{\mathrm n}+Z m_{\mathrm p}+Z m_{\mathrm e}
-\left(
\aV A-\aS A^{2/3}-\aC Z^2 A^{-1/3}
-\aA\frac{(A-2Z)^2}{A}+\delta\aP A^{-1/2}
\right)c^{-2}.
\]
\end{formulabox}

\begin{notebox}[title=Vorzeichenlogik]
Die Bindungsenergie $B(A,Z)$ wird von den freien Nukleonenmassen abgezogen. Terme, die die Bindung verringern, stehen in $B$ mit Minuszeichen, erscheinen aber in $M c^2$ mit Pluszeichen. Deshalb erhoehen Coulombterm und Asymmetrieterm die Masse.
\end{notebox}

\newpage

%%%%%%%%%%%%%%%%%%%%%%%%%%%%%%%%%%%%%%%%%%%%%%%%%%%%%%%%%%%%%%%%%%%
\section{Aufgabe 1: Bethe-Weizsaecker - der Coulombterm}

\begin{taskbox}
Fuer schwere Atomkerne modelliert man den Kern als homogen geladene Kugel mit Gesamtladung $Q$ und Radius $R$.
\begin{enumerate}[label=(\alph*)]
    \item Bestimme die elektrostatische Abstossungsenergie $\Eab$ einer homogen geladenen Kugel.
    \item Setze diese Energie gleich dem Coulombterm der Bethe-Weizsaecker-Formel und verwende $R=R_0A^{1/3}$, um $\aC$ herzuleiten.
    \item Berechne $\aC$ fuer $R_0=\SI{1.2}{\fermi}$ und $R_0=\SI{1.3}{\fermi}$ und bewerte das Modell.
\end{enumerate}
\end{taskbox}

\subsection{Abgabe / vollstaendige Loesung}

\begin{enumerate}[label=(\alph*)]
    \item \textbf{Selbstenergie der homogen geladenen Kugel.}

    \begin{calculationbox}
    Wir bauen die Kugel gedanklich aus duennen Kugelschalen auf. Hat die bereits aufgebaute Kugel den Radius $r$, so ist wegen der homogenen Ladungsdichte die eingeschlossene Ladung proportional zum Volumen:
    \[
        q(r)=Q\frac{r^3}{R^3}.
    \]
    Das Potential an der Oberflaeche dieser Zwischenkugel ist
    \[
        \phi(r)=\frac{1}{4\pi\varepsilon_0}\frac{q(r)}{r}.
    \]
    Die Ladungszunahme beim Vergroessern von $r$ um $\dd r$ ist
    \[
        \dd q = \dv{q}{r}\dd r = \frac{3Q}{R^3}r^2\dd r.
    \]
    Daher ist die aufzubringende Arbeit
    \[
        \dd E = \phi(r)\dd q
        =\frac{1}{4\pi\varepsilon_0}\frac{Qr^3/R^3}{r}\frac{3Q}{R^3}r^2\dd r
        =\frac{1}{4\pi\varepsilon_0}\frac{3Q^2}{R^6}r^4\dd r.
    \]
    Integration von $0$ bis $R$ ergibt
    \[
        \Eab=\int_0^R \dd E
        =\frac{3Q^2}{4\pi\varepsilon_0R^6}\int_0^R r^4\dd r
        =\frac{3Q^2}{4\pi\varepsilon_0R^6}\frac{R^5}{5}.
    \]
    \end{calculationbox}

    \begin{answerbox}
    \[
        \boxed{\Eab=\frac{3}{5}\frac{1}{4\pi\varepsilon_0}\frac{Q^2}{R}}.
    \]
    \end{answerbox}

    \item \textbf{Vergleich mit dem Coulombterm.}

    \begin{calculationbox}
    Fuer einen Kern mit Protonenzahl $Z$ gilt $Q=Ze$. Damit folgt
    \[
        \Eab=\frac{3}{5}\frac{1}{4\pi\varepsilon_0}\frac{Z^2e^2}{R}.
    \]
    Mit dem Rutherford-Radius
    \[
        R=R_0A^{1/3}
    \]
    wird
    \[
        \Eab=\frac{3}{5}\frac{e^2}{4\pi\varepsilon_0R_0}\frac{Z^2}{A^{1/3}}.
    \]
    Der Coulombterm der Bethe-Weizsaecker-Formel hat die Form
    \[
        E_C=\aC\frac{Z^2}{A^{1/3}}.
    \]
    Koeffizientenvergleich liefert direkt
    \[
        \aC=\frac{3}{5}\frac{e^2}{4\pi\varepsilon_0R_0}.
    \]
    \end{calculationbox}

    \begin{answerbox}
    \[
        \boxed{\aC=\frac{3}{5}\frac{e^2}{4\pi\varepsilon_0R_0}}.
    \]
    \end{answerbox}

    \item \textbf{Numerische Werte.}

    \begin{calculationbox}
    Mit
    \[
        \frac{e^2}{4\pi\varepsilon_0}\approx \SI{1.44}{\MeV\fermi}
    \]
    folgt
    \[
        \aC=\frac{3}{5}\frac{\SI{1.44}{\MeV\fermi}}{R_0}.
    \]
    Fuer $R_0=\SI{1.2}{\fermi}$:
    \[
        \aC=\frac{3}{5}\frac{1.44}{1.2}\,\si{\MeV}=\SI{0.72}{\MeV}.
    \]
    Fuer $R_0=\SI{1.3}{\fermi}$:
    \[
        \aC=\frac{3}{5}\frac{1.44}{1.3}\,\si{\MeV}\approx\SI{0.665}{\MeV}.
    \]
    \end{calculationbox}

    \begin{answerbox}
    \[
        \boxed{\aC(R_0=\SI{1.2}{\fermi})\approx\SI{0.72}{\MeV}},\qquad
        \boxed{\aC(R_0=\SI{1.3}{\fermi})\approx\SI{0.665}{\MeV}}.
    \]
    Das Modell ist erstaunlich gut: Der empirische Wert liegt in der Groessenordnung $\aC\approx\SI{0.71}{\MeV}$. Die homogen geladene Kugel beschreibt also den Coulombanteil schwerer Kerne plausibel, obwohl reale Kerne keine scharf begrenzten homogenen Ladungskugeln sind.
    \end{answerbox}
\end{enumerate}

\subsection{Tafelvortrag}
\begin{summarybox}[title=5-Minuten-Struktur]
\begin{enumerate}
    \item Skizze: Kugel mit Radius $R$, Zwischenkugel $r$, Gesamtladung $Q$.
    \item Kernidee: Kugel schalenweise aufbauen, $\dd E=\phi\,\dd q$.
    \item Homogene Dichte: $q(r)=Q r^3/R^3$, also $\dd q=3Qr^2\dd r/R^3$.
    \item Integration: $\Eab=\int_0^R \frac{1}{4\pi\varepsilon_0}\frac{q(r)}{r}\dd q$.
    \item Vergleich mit $E_C=\aC Z^2/A^{1/3}$ und Ergebnis $\aC=\frac{3}{5}e^2/(4\pi\varepsilon_0R_0)$.
\end{enumerate}
\end{summarybox}

\begin{notebox}[title=Typische Rueckfrage]
Warum steht in der Massenformel oft $Z(Z-1)$ statt $Z^2$? Physikalisch wechselwirkt ein Proton nicht mit sich selbst. Fuer grosse $Z$ ist $Z(Z-1)\approx Z^2$. Das Kugelmodell ist ohnehin eine kontinuierliche Naeherung.
\end{notebox}

\newpage

%%%%%%%%%%%%%%%%%%%%%%%%%%%%%%%%%%%%%%%%%%%%%%%%%%%%%%%%%%%%%%%%%%%
\section{Aufgabe 2: Bethe-Weizsaecker - der Volumenterm}

\begin{taskbox}
Der Volumenterm der Bethe-Weizsaecker-Formel lautet $\aV A$. Er soll mit einem Nachbarbindungsmodell verstanden werden. Wegen der kurzen Reichweite der Kernkraft wechselwirken nur naechste Nachbarn. Jede Bindung zweier Nukleonen setzt eine Energie $\varepsilon$ frei.
\begin{enumerate}[label=(\alph*)]
    \item Wie viele naechste Nachbarn besitzt ein Nukleon in dichtester Kugelpackung?
    \item Bestimme $k$ in $\EWW=k\varepsilon A$.
    \item Verwende die Bindungsenergie von Deuterium, $\SI{2.23}{\MeV}$, zur Abschaetzung von $\varepsilon$ und $\aV$.
\end{enumerate}
\end{taskbox}

\subsection{Abgabe / vollstaendige Loesung}

\begin{enumerate}[label=(\alph*)]
    \item \textbf{Naechste Nachbarn.}

    \begin{solutionbox}
    In der dichtesten Kugelpackung beruehrt jede Kugel $12$ andere Kugeln. Uebertragen auf das Nukleonenmodell bedeutet das:
    \[
        N_{\mathrm{NN}}=12.
    \]
    Oberflaecheneffekte werden in der Aufgabenstellung bewusst ignoriert.
    \end{solutionbox}

    \begin{answerbox}
    \[
        \boxed{N_{\mathrm{NN}}=12}.
    \]
    \end{answerbox}

    \item \textbf{Doppelzaehlung der Bindungen.}

    \begin{calculationbox}
    Ein Kern mit $A$ Nukleonen haette bei naiver Zaehlung $12A$ Nachbarschaften. Jede reale Bindung wird dabei aber zweimal gezaehlt: einmal von Nukleon 1 zu Nukleon 2 und einmal von Nukleon 2 zu Nukleon 1. Deshalb ist die Zahl der Bindungen
    \[
        N_{\mathrm{Bindungen}}=\frac{12A}{2}=6A.
    \]
    Wenn jede Bindung die Energie $\varepsilon$ freisetzt, gilt
    \[
        \EWW=6A\varepsilon.
    \]
    Vergleich mit $\EWW=k\varepsilon A$ liefert
    \[
        k=6.
    \]
    \end{calculationbox}

    \begin{answerbox}
    \[
        \boxed{k=6}.
    \]
    \end{answerbox}

    \item \textbf{Abschaetzung mit Deuterium.}

    \begin{calculationbox}
    Das Deuteron $\isotope{2}{H}$ besteht aus einem Proton und einem Neutron. Es besitzt in diesem Minimalmodell genau eine Bindung. Daher setzen wir
    \[
        \varepsilon \approx B(\isotope{2}{H})=\SI{2.23}{\MeV}.
    \]
    Aus Teil (b) folgt fuer grosse Kerne
    \[
        \EWW=6\varepsilon A.
    \]
    Der Volumenterm lautet
    \[
        E_V=\aV A.
    \]
    Also identifiziert man
    \[
        \aV=6\varepsilon=6\cdot\SI{2.23}{\MeV}=\SI{13.38}{\MeV}.
    \]
    Vergleich mit $\aV=\SI{15.85}{\MeV}$:
    \[
        \Delta \aV=\SI{15.85}{\MeV}-\SI{13.38}{\MeV}=\SI{2.47}{\MeV},
    \]
    \[
        \frac{\Delta \aV}{\SI{15.85}{\MeV}}\approx 0.156\approx \SI{15.6}{\percent}.
    \]
    \end{calculationbox}

    \begin{answerbox}
    \[
        \boxed{\varepsilon\approx\SI{2.23}{\MeV}},\qquad
        \boxed{\aV\approx\SI{13.38}{\MeV}}.
    \]
    Das Ergebnis ist kleiner als der verwendete empirische Wert $\SI{15.85}{\MeV}$, aber die Groessenordnung stimmt gut. Die Abweichung ist erwartbar, weil Deuterium ein sehr spezielles Zwei-Nukleonen-System ist und das Modell Schalenstruktur, Spinabhaengigkeit, Tensoranteile und Oberflaecheneffekte nicht enthaelt.
    \end{answerbox}
\end{enumerate}

\subsection{Tafelvortrag}
\begin{summarybox}[title=Kernaussage fuer die Tafel]
\[
    \text{12 Nachbarn pro Nukleon}\quad\Rightarrow\quad \frac{12A}{2}=6A\text{ Bindungen}\quad\Rightarrow\quad \EWW=6\varepsilon A.
\]
Die Division durch $2$ ist der wichtigste Schritt: Ohne sie wuerde jede Nukleon-Nukleon-Bindung doppelt gezaehlt.
\end{summarybox}

\begin{notebox}[title=Lernintuition]
Der Volumenterm ist proportional zu $A$, weil jedes Nukleon im Inneren ungefaehr gleich viele Nachbarn hat. Die Kernkraft ist kurzreichweitig; deshalb waechst die Bindung nicht mit der Zahl aller Paare $\sim A^2$, sondern nur mit der Zahl lokaler Nachbarbindungen $\sim A$.
\end{notebox}

\newpage

%%%%%%%%%%%%%%%%%%%%%%%%%%%%%%%%%%%%%%%%%%%%%%%%%%%%%%%%%%%%%%%%%%%
\section{Aufgabe 3: Stabilstes Nuklid einer Isobare}

\begin{taskbox}
Fuer eine Isobarenreihe wird $A$ konstant gehalten. Mit der Bethe-Weizsaecker-Formel soll die Kernladungszahl $Z_0$ bestimmt werden, bei der die Masse minimal wird. Setze dafuer $\delta=0$.
\begin{enumerate}[label=(\alph*)]
    \item Leite $Z_0$ allgemein aus der Massenformel her.
    \item Berechne $Z_0$ fuer $A=25$ und $A=67$, vergleiche mit der Nuklidkarte und zeichne die Massenparabeln im Bereich $Z_0\pm 3$ mit theoretischen und gemessenen Massen.
\end{enumerate}
\end{taskbox}

\subsection{Abgabe / vollstaendige Loesung}

\begin{enumerate}[label=(\alph*)]
    \item \textbf{Allgemeine Herleitung von $Z_0$.}

    \begin{calculationbox}
    Wir minimieren $M(A,Z)$ fuer festes $A$. Es ist uebersichtlicher, mit $M(A,Z)c^2$ zu arbeiten. Mit $\delta=0$ gilt
    \[
    M(A,Z)c^2=(A-Z)m_{\mathrm n}c^2+Z m_{\mathrm p}c^2+Z m_{\mathrm e}c^2
    -\aV A+\aS A^{2/3}+\aC Z^2 A^{-1/3}+\aA\frac{(A-2Z)^2}{A}.
    \]
    Fuer die Ableitung nach $Z$ sind nur die $Z$-abhaengigen Terme relevant:
    \[
        f(Z)=Z(m_{\mathrm p}+m_{\mathrm e}-m_{\mathrm n})c^2
        +\aC Z^2A^{-1/3}+\aA\frac{(A-2Z)^2}{A}.
    \]
    Ableiten:
    \[
        \dv{f}{Z}=(m_{\mathrm p}+m_{\mathrm e}-m_{\mathrm n})c^2
        +2\aC ZA^{-1/3}-4\aA\frac{A-2Z}{A}.
    \]
    Minimumsbedingung:
    \[
        \dv{f}{Z}=0.
    \]
    Also
    \[
        (m_{\mathrm p}+m_{\mathrm e}-m_{\mathrm n})c^2
        +2\aC ZA^{-1/3}-4\aA+\frac{8\aA}{A}Z=0.
    \]
    Wir sammeln die $Z$-Terme:
    \[
        Z\left(2\aC A^{-1/3}+\frac{8\aA}{A}\right)
        =4\aA-(m_{\mathrm p}+m_{\mathrm e}-m_{\mathrm n})c^2.
    \]
    Damit
    \[
        Z_0=\frac{4\aA-(m_{\mathrm p}+m_{\mathrm e}-m_{\mathrm n})c^2}
        {2\aC A^{-1/3}+8\aA/A}.
    \]
    Multiplikation von Zaehler und Nenner mit $A$ liefert die aequivalente Form
    \[
        Z_0=\frac{A\left[4\aA-(m_{\mathrm p}+m_{\mathrm e}-m_{\mathrm n})c^2\right]}
        {2\aC A^{2/3}+8\aA}.
    \]
    \end{calculationbox}

    \begin{answerbox}
    \[
        \boxed{Z_0=\frac{4\aA-(m_{\mathrm p}+m_{\mathrm e}-m_{\mathrm n})c^2}
        {2\aC A^{-1/3}+8\aA/A}}
    \]
    beziehungsweise
    \[
        \boxed{Z_0=\frac{A\left[4\aA-(m_{\mathrm p}+m_{\mathrm e}-m_{\mathrm n})c^2\right]}
        {2\aC A^{2/3}+8\aA}}.
    \]
    \end{answerbox}

    \item \textbf{Numerische Auswertung fuer $A=25$ und $A=67$.}

    \begin{calculationbox}
    Verwendete Ruheenergien:
    \[
        m_{\mathrm p}c^2=\SI{938.272}{\MeV},\qquad
        m_{\mathrm n}c^2=\SI{939.565}{\MeV},\qquad
        m_{\mathrm e}c^2=\SI{0.511}{\MeV}.
    \]
    Daher
    \[
        (m_{\mathrm p}+m_{\mathrm e}-m_{\mathrm n})c^2
        \approx 938.272+0.511-939.565
        =\SI{-0.782}{\MeV}.
    \]
    Mit $\aA=\SI{23.21}{\MeV}$ und $\aC=\SI{0.71}{\MeV}$ wird
    \[
        Z_0=\frac{4\cdot23.21+0.782}{2\cdot0.71 A^{-1/3}+8\cdot23.21/A}.
    \]
    Fuer $A=25$:
    \[
        A^{-1/3}=25^{-1/3}\approx0.342,
    \]
    \[
        Z_0(25)=\frac{93.622}{2\cdot0.71\cdot0.342+185.68/25}
        \approx\frac{93.622}{7.913}\approx11.83.
    \]
    Fuer $A=67$:
    \[
        A^{-1/3}=67^{-1/3}\approx0.2465,
    \]
    \[
        Z_0(67)=\frac{93.622}{2\cdot0.71\cdot0.2465+185.68/67}
        \approx\frac{93.622}{3.121}\approx30.00.
    \]
    \end{calculationbox}

    \begin{answerbox}
    \[
        \boxed{Z_0(A=25)\approx11.83\Rightarrow Z=12\Rightarrow \isotope{25}{Mg}},
    \]
    \[
        \boxed{Z_0(A=67)\approx30.00\Rightarrow Z=30\Rightarrow \isotope{67}{Zn}}.
    \]
    Beide Ergebnisse passen zur Nuklidkarte: $\isotope{25}{Mg}$ und $\isotope{67}{Zn}$ sind stabile Nuklide. Die Massenformel beschreibt also die Lage des Massental-Minimums gut.
    \end{answerbox}
\end{enumerate}

\subsection{Tabellen fuer die Massenparabeln}

\begin{notebox}[title=Darstellung der Plots]
Die absoluten Atommassen liegen bei vielen zehntausend $\si{\MeV}/c^2$, waehrend sich die Isobaren innerhalb einer Reihe nur um wenige bis einige zehn $\si{\MeV}/c^2$ unterscheiden. Deshalb werden in den Plots relative Massen
\[
    \Delta M(Z)=M(A,Z)-M_{\min}(A)
\]
gezeigt. Das verschiebt die Kurven nur vertikal und aendert das Minimum nicht.
\end{notebox}

\begin{table}[h]
\centering
\caption{Theoretische und experimentelle relative Massen fuer $A=25$.}
\begin{tabular}{c c c c c}
\toprule
$Z$ & Nuklid & $M_{\mathrm{theo}}$ / $\si{\MeV}/c^2$ & $\Delta M_{\mathrm{theo}}$ / $\si{\MeV}/c^2$ & $\Delta M_{\mathrm{exp}}$ / $\si{\MeV}/c^2$ \\
\midrule
9  & $\isotope{25}{F}$  & 23302.333 & 31.613 & 24.530 \\
10 & $\isotope{25}{Ne}$ & 23283.882 & 13.162 & 11.154 \\
11 & $\isotope{25}{Na}$ & 23273.345 &  2.625 &  3.836 \\
12 & $\isotope{25}{Mg}$ & 23270.720 &  0.000 &  0.000 \\
13 & $\isotope{25}{Al}$ & 23276.008 &  5.288 &  4.277 \\
14 & $\isotope{25}{Si}$ & 23289.209 & 18.489 & 17.021 \\
15 & $\isotope{25}{P}$  & 23310.323 & 39.603 & 33.388 \\
\bottomrule
\end{tabular}
\end{table}

\begin{table}[h]
\centering
\caption{Theoretische und experimentelle relative Massen fuer $A=67$.}
\begin{tabular}{c c c c c}
\toprule
$Z$ & Nuklid & $M_{\mathrm{theo}}$ / $\si{\MeV}/c^2$ & $\Delta M_{\mathrm{theo}}$ / $\si{\MeV}/c^2$ & $\Delta M_{\mathrm{exp}}$ / $\si{\MeV}/c^2$ \\
\midrule
27 & $\isotope{67}{Co}$ & 62350.832 & 14.025 & 12.559 \\
28 & $\isotope{67}{Ni}$ & 62343.036 &  6.229 &  4.137 \\
29 & $\isotope{67}{Cu}$ & 62338.361 &  1.554 &  0.561 \\
30 & $\isotope{67}{Zn}$ & 62336.807 &  0.000 &  0.000 \\
31 & $\isotope{67}{Ga}$ & 62338.374 &  1.567 &  1.001 \\
32 & $\isotope{67}{Ge}$ & 62343.062 &  6.255 &  5.207 \\
33 & $\isotope{67}{As}$ & 62350.871 & 14.064 & 11.293 \\
\bottomrule
\end{tabular}
\end{table}

\FloatBarrier

\subsection{Massenparabeln}

\begin{figure}[h]
\centering
\begin{tikzpicture}
\begin{axis}[
    width=0.88\textwidth,
    height=7cm,
    xlabel={$Z$},
    ylabel={$\Delta M$ / $\si{\MeV}/c^2$},
    title={Isobarenparabel fuer $A=25$},
    grid=both,
    xmin=8.5,xmax=15.5,
    ymin=0,ymax=42,
    legend pos=north west
]
\addplot+[mark=*] table[row sep=\\]{%
Z y\\
9 31.613\\
10 13.162\\
11 2.625\\
12 0.000\\
13 5.288\\
14 18.489\\
15 39.603\\
};
\addlegendentry{Massenformel}
\addplot+[only marks,mark=square*] table[row sep=\\]{%
Z y\\
9 24.530\\
10 11.154\\
11 3.836\\
12 0.000\\
13 4.277\\
14 17.021\\
15 33.388\\
};
\addlegendentry{gemessen}
\end{axis}
\end{tikzpicture}
\caption{Relative Massen der Isobarenreihe $A=25$. Das Minimum liegt bei $Z=12$, also bei $\isotope{25}{Mg}$.}
\end{figure}

\begin{figure}[h]
\centering
\begin{tikzpicture}
\begin{axis}[
    width=0.88\textwidth,
    height=7cm,
    xlabel={$Z$},
    ylabel={$\Delta M$ / $\si{\MeV}/c^2$},
    title={Isobarenparabel fuer $A=67$},
    grid=both,
    xmin=26.5,xmax=33.5,
    ymin=0,ymax=16,
    legend pos=north west
]
\addplot+[mark=*] table[row sep=\\]{%
Z y\\
27 14.025\\
28 6.229\\
29 1.554\\
30 0.000\\
31 1.567\\
32 6.255\\
33 14.064\\
};
\addlegendentry{Massenformel}
\addplot+[only marks,mark=square*] table[row sep=\\]{%
Z y\\
27 12.559\\
28 4.137\\
29 0.561\\
30 0.000\\
31 1.001\\
32 5.207\\
33 11.293\\
};
\addlegendentry{gemessen}
\end{axis}
\end{tikzpicture}
\caption{Relative Massen der Isobarenreihe $A=67$. Das Minimum liegt bei $Z=30$, also bei $\isotope{67}{Zn}$.}
\end{figure}

\FloatBarrier

\subsection{Diskussion}
\begin{solutionbox}[title=Physikalische Interpretation]
Fuer festes $A$ konkurrieren im Wesentlichen zwei $Z$-abhaengige Effekte:
\begin{itemize}
    \item Der \textbf{Coulombterm} wird bei grossem $Z$ groesser und bevorzugt daher kleinere Protonenzahlen.
    \item Der \textbf{Asymmetrieterm} bestraft grosse Abweichungen von $N=Z$, also von $A-2Z=0$.
\end{itemize}
Das Minimum der Isobarenparabel liegt dort, wo diese Effekte zusammen mit der kleinen Massendifferenz zwischen Neutron und Proton ausgeglichen sind. Fuer leichte Kerne liegt das Minimum naeher bei $N\approx Z$. Fuer schwerere Kerne wird wegen der Coulombabstossung ein Neutronenueberschuss energetisch guenstig.
\end{solutionbox}

\begin{warningbox}[title=Warum die gemessenen Punkte nicht perfekt auf der Parabel liegen]
Die Bethe-Weizsaecker-Formel ist eine makroskopische Naeherung. Reale Massen enthalten zusaetzlich Schalenkorrekturen, Paarungseffekte, Deformationen und weitere mikroskopische Beitraege. Deshalb sagt die Massenformel den Trend und das ungefaehre Minimum sehr gut voraus, aber nicht jeden einzelnen Massenwert exakt.
\end{warningbox}

\subsection{Tafelvortrag}
\begin{summarybox}[title=8-Minuten-Struktur]
\begin{enumerate}
    \item Schreibe die Bethe-Weizsaecker-Formel hin und setze $\delta=0$.
    \item Erklaere: $A$ konstant, deshalb nur $Z$-abhaengige Terme betrachten.
    \item Definiere
    \[
        f(Z)=Z(m_p+m_e-m_n)c^2+\aC Z^2A^{-1/3}+\aA\frac{(A-2Z)^2}{A}.
    \]
    \item Ableiten und Null setzen:
    \[
        f'(Z)=0.
    \]
    \item Ergebnis fuer $Z_0$ boxed anschreiben.
    \item Zahlen einsetzen: $Z_0(25)\approx11.83$, $Z_0(67)\approx30.00$.
    \item Nuklidkarte: $\isotope{25}{Mg}$ und $\isotope{67}{Zn}$ sind stabile Minimumskandidaten.
    \item Plot zeigen: Theorie ist Parabel; gemessene Massen folgen der Parabel naeherungsweise.
\end{enumerate}
\end{summarybox}

\newpage

%%%%%%%%%%%%%%%%%%%%%%%%%%%%%%%%%%%%%%%%%%%%%%%%%%%%%%%%%%%%%%%%%%%
\section*{Pruefungs- und Lernfragen}
\addcontentsline{toc}{section}{Pruefungs- und Lernfragen}

\begin{enumerate}
    \item Warum skaliert der Coulombterm wie $Z^2/A^{1/3}$?
    \item Woher kommt der Faktor $3/5$ in der Selbstenergie der homogen geladenen Kugel?
    \item Warum ist der Volumenterm proportional zu $A$ und nicht zu $A^2$?
    \item Warum muss man bei der Zahl der Nachbarbindungen durch $2$ teilen?
    \item Warum ist Deuterium als Abschaetzung fuer $\varepsilon$ plausibel, aber nicht perfekt?
    \item Welche Terme der Bethe-Weizsaecker-Formel sind fuer die Minimierung nach $Z$ bei festem $A$ relevant?
    \item Warum ist die Isobarenmasse als Funktion von $Z$ naeherungsweise eine Parabel?
    \item Warum verschiebt sich das stabile $Z$ bei schweren Kernen zu kleinerem $Z/A$, also zu Neutronenueberschuss?
\end{enumerate}

\begin{summarybox}[title=Endergebnisse]
\[
    \Eab=\frac{3}{5}\frac{1}{4\pi\varepsilon_0}\frac{Q^2}{R},
    \qquad
    \aC=\frac{3}{5}\frac{e^2}{4\pi\varepsilon_0R_0}.
\]
\[
    \aC(\SI{1.2}{\fermi})\approx\SI{0.72}{\MeV},
    \qquad
    \aC(\SI{1.3}{\fermi})\approx\SI{0.665}{\MeV}.
\]
\[
    N_{\mathrm{NN}}=12,
    \qquad
    k=6,
    \qquad
    \varepsilon\approx\SI{2.23}{\MeV},
    \qquad
    \aV\approx\SI{13.38}{\MeV}.
\]
\[
    Z_0=\frac{4\aA-(m_{\mathrm p}+m_{\mathrm e}-m_{\mathrm n})c^2}
    {2\aC A^{-1/3}+8\aA/A}.
\]
\[
    Z_0(25)\approx11.83\Rightarrow \isotope{25}{Mg},
    \qquad
    Z_0(67)\approx30.00\Rightarrow \isotope{67}{Zn}.
\]
\end{summarybox}

\end{document}
